\section[$d=1$ Conformal Galilei algebras ${\mathfrak g}_{\ell}$]{$\boldsymbol{d=1}$ Conformal Galilei algebras $\boldsymbol{{\mathfrak g}_{\ell}}$}

The complex Lie algebra $ {\mathfrak g}_{\ell} $ for a~f\/ixed integer $ \ell $ has the elements~\cite{NdORM1}:
\begin{gather*}
D, \ H, \ C, \ P_n, \quad n = 0, 1,   \dots,   2\ell.
%\label{generators}
\end{gather*}
Their nonvanishing commutators are given~by
\begin{alignat}{4}
& [D, H] = H, \qquad & & [D, C] = - C, \qquad & & [C, H] = 2D,& \nonumber\\
& [H, P_n] = -n P_{n-1}, \qquad & & [D, P_n] = (\ell-n) P_n, \qquad & & [C, P_n] = (2\ell-n) P_{n+1}.&\label{nonvanishing-commutators}
\end{alignat}
One may see from this that $ \langle   P_0, P_1, \dots, P_{2\ell}   \rangle $ is an Abelian ideal of $ {\mathfrak
g}_{\ell} $, so that $ {\mathfrak g}_{\ell} $ is \textit{not} semisimple.
It is known that this Lie algebra has no central extensions~\cite{MT}.
The subalgebra spanned by $ \langle
 H, D, C
\rangle $ is isomorphic to $ {\mathfrak{so}}(2,1) \simeq {\mathfrak{sl}}(2,{\mathbb R}) \simeq {\mathfrak{su}}(1,1)$.
The Abelian subalgebra spanned by $ \langle
P_n
\rangle_{n=0, 1, \dots, 2\ell} $ carries the spin $ \ell $ representation of the $ {\mathfrak{sl}}(2,{\mathbb R})$
subalgebra.

This algebra may be realized as generators of transformation of $(1+1)$-dimensional spacetime
\begin{gather*}
H = \frac{\partial}{\partial t},
\qquad
D =-t\frac{\partial}{\partial t} -\ell x \frac{\partial}{\partial x},
\qquad
C = t^2 \frac{\partial}{\partial t} + 2 \ell tx \frac{\partial}{\partial x},
\qquad
P_n = (-t)^n \frac{\partial}{\partial x}.
%\label{CGArealization-in-space}
\end{gather*}
In this realization,~$H$,~$D$ and~$C$ generates time translation, dilatation and the special conformal representation,
respectively.
Meanwhile the generator of space translation is represented by~$P_0$, the Galilei transformation generator by~$P_1$,
the transformation to a~reference frame with constant acceleration is given by~$P_2$ and so on.

One may introduce the algebraic anti-involution $ \omega: {\mathfrak g}_{\ell} \to {\mathfrak g}_{\ell} $~by
\begin{gather*}
\omega(D) = D, \qquad \omega(H) = C, \qquad \omega(P_n) = P_{2\ell-n}.
%\label{anti-inv}
\end{gather*}
It is not dif\/f\/icult to verify that $ \omega $ satisf\/ies the required relations:
\begin{gather*}
\omega([X,Y]) = [\omega(Y), \omega(X)], \qquad \omega^2(X) = X \qquad  \forall\, X \in {\mathfrak g}_{\ell}.
%\label{anti-inv2}
\end{gather*}
Let us def\/ine the degree of the generators based on their commutator with respect to~$D$:
\begin{gather*}
\deg(D) = 0, \qquad \deg(H) = 1, \qquad \deg(C) = -1, \qquad \deg(P_n) = \ell-n.
%\label{degree}
\end{gather*}
With respect to the sign of the degree one may def\/ine the triangular decomposition of ${\mathfrak g}_{\ell}$:
\begin{gather*}
    {\mathfrak g}_{\ell} = {\mathfrak g}_{\ell}^+ \oplus {\mathfrak g}_{\ell}^0 \oplus {\mathfrak g}_{\ell}^-,
%\label{tri-decomp}
\end{gather*}
where
\begin{gather*}
{\mathfrak g}_{\ell}^+ = \langle H, P_0, P_1, \dots, P_{\ell-1} \rangle,
\qquad
{\mathfrak g}_{\ell}^0 = \langle D, P_{\ell} \rangle,
\qquad
{\mathfrak g}_{\ell}^- = \langle C, P_{\ell+1}, P_{\ell+2}, \dots, P_{2\ell} \rangle.
\end{gather*}
This is a~decomposition of $ {\mathfrak g}_{\ell} $ as a~direct sum of the vector spaces.


\section{Verma modules and Shapovalov determinant}
\label{SEC:VM}

We start with the one-dimensional module $ {\mathbb C} \ket{\delta,p} $ over $ {\mathfrak b} = {\mathfrak g}_{\ell}^0
\oplus {\mathfrak g}_{\ell}^-$. The lowest weight vector $ \ket{\delta,p} $ is def\/ined~by
\begin{gather*}
D \ket{\delta,p} = \delta \ket{\delta,p},
\qquad
P_{\ell} \ket{\delta,p} = p \ket{\delta,p},
\qquad
X \ket{\delta,p} = 0,
\qquad
 \forall \, X \in {\mathfrak g}_{\ell}^-.
%\label{LWV}
\end{gather*}
Then for each pair of $(\delta,p) $ the Verma module over $ {\mathfrak g}_{\ell} $ is def\/ined as usual (see, e.g.,~\cite{Dix}): $ V^{\delta,p}_{\ell} = U({\mathfrak g}_{\ell}) \otimes_{U({\mathfrak b})} \ket{\delta,p} $ where $U({\mathfrak g}_{\ell})$ and $ U({\mathfrak b}) $ are the enveloping algebras of $ {\mathfrak g}_{\ell} $ and $
\mathfrak b$, respectively.
In order to specify the basis of $ V^{\delta,p}_{\ell} $ we introduce the $ \ell$-component vector
\begin{gather*}
\underaccent{\tilde}{m} = (m_1, m_2, \dots, m_{\ell})  \in {\mathbb R}^{\ell},
%\label{m-vec-def}
\end{gather*}
where $ m_i$ $(i=1,2,\dots,\ell) $ are non-negative integers.
With this the basis of $ V^{\delta,p}_{\ell} $ is given~by
\begin{gather}
\ket{k,\underaccent{\tilde}{m}} = H^k P_{\ell-1}^{m_1} P_{\ell-2}^{m_2} \cdots P_0^{m_{\ell}} \ket{\delta,p}.
\label{anyl-basis}
\end{gather}
We also introduce the~$\ell$-component vectors
\begin{gather*}
\underaccent{\tilde}{\epsilon}_j = (0, \dots, 0, 1, 0, \dots,0) \in {\mathbb R}^{\ell},
\qquad
j=1,2,\dots, \ell.
%\label{tilde-epsilon-def}
\end{gather*}
where the~$j$th entry of $\underaccent{\tilde}{\epsilon}_j $ is~1, and all other entries are~0.
It may not be dif\/f\/icult to prove the following relation by induction on $k$:
\begin{gather}
[P_n, H^k] = \sum\limits_{i=1}^{\min\{n,k\}} \binom{n}{i} \frac{k!}{(k-i)!}  H^{k-i}   P_{n-i},
\qquad
n \geq 1.
\label{PnHkcomm}
\end{gather}
It follows that the action of $ {\mathfrak g}_{\ell}^0$ and ${\mathfrak g}_{\ell}^+$ on $\ket{k,\underaccent{\tilde}{m}}$:
\begin{gather*}
    D \ket{k,\underaccent{\tilde}{m}} = \left(\delta + k+ \sum\limits_{i=1}^{\ell} i   m_i\right) \ket{k,\underaccent{\tilde}{m}},
\nonumber
\\
    P_{\ell} \ket{k,\underaccent{\tilde}{m}} = p \ket{k,\underaccent{\tilde}{m}} + \sum\limits_{i=1}^{\min\{\ell,k\}}
\binom{\ell}{i} \frac{k!}{(k-i)!} \ket{k-i,\underaccent{\tilde}{m}+\underaccent{\tilde}{\epsilon}_{i}},
\nonumber
\\
    H \ket{k,\underaccent{\tilde}{m}}= \ket{k+1,\underaccent{\tilde}{m}},
\nonumber
\\
    P_n \ket{k,\underaccent{\tilde}{m}} = \sum\limits_{i=1}^{\min\{n,k\}} \binom{n}{i} \frac{k!}{(k-i)!}
\ket{k-i,\underaccent{\tilde}{m}+\underaccent{\tilde}{\epsilon}_{\ell-n+i}},
\qquad
1 \leq n \leq \ell-1,
\nonumber
\\
    P_0 \ket{k,\underaccent{\tilde}{m}} = \ket{k,\underaccent{\tilde}{m}+\underaccent{\tilde}{\epsilon}_{\ell}}.
%\label{g0p-action-any}
\end{gather*}
From this one see that $ V^{\delta,p}_{\ell} $ has a~grading structure according to the eigenvalue of~$D$:
\begin{gather*}
    V^{\delta,p}_{\ell} = \bigoplus_{N=0}^{\infty} \big(V^{\delta,p}_{\ell}\big)_N,
\qquad
\big(V^{\delta,p}_{\ell}\big)_N = \big\{ \ket{v} \in V^{\delta,p}_{\ell} \, | \, D \ket{v} = (\delta+N) \ket{v} \big\}.
%\label{Verma-grading}
\\
    N = k+ \sum\limits_{i=1}^{\ell} i   m_i.
%\label{N-definition}
\end{gather*}
We refer to~$N$ as the \textit{level} as usual.

Now we def\/ine Shapovalov form on a~module over $ {\mathfrak g}_{\ell} $~\cite{Shapo} (see also, e.g.,~\cite{KacRai}).
That is a~contravariant Hermitian form $ \bracket{\cdot}{\cdot} $ def\/ined by making use of the anti-involution~ $\omega$.
For the Verma module $ V^{\delta,p}_{\ell} $ it is def\/ined as follows:   Let $ \ket{x} $ and $ \ket{y} $ be any two
vectors in $ V^{\delta,p}_{\ell}$.
They may be written as
\begin{gather*}
\ket{x} = X \ket{\delta,p},
\qquad
\ket{y} = Y \ket{\delta,p},
\qquad
X, Y \in U({\mathfrak g}_{\ell}^+).
\end{gather*}
Then
\begin{gather}
\bracket{x}{y} = \bra{\delta,p} \omega(X) Y \ket{\delta,p},
\qquad
\bracket{\delta,p}{\delta,p} = 1.
\label{innerprod-def}
\end{gather}
Next we def\/ine the Shapovalov determinant at level~$N$.
Let $ \ket{v_1}, \ket{v_2}, \dots, \ket{v_r} $ be a~set of basis of the subspace $ \big(V^{\delta,p}_{\ell}\big)_N$.
We consider the matrix $ (\bracket{v_i}{v_j}) $ whose entries are     the Shapovalov form and we call the determinant of
this matrix, the Shapovalov determinant at level $N: \Delta^{(\ell)}_N = \det (\bracket{v_i}{v_j})$.
We give an explicit formulae of the Shapovalov determinants of $ {\mathfrak g}_{\ell}$, since $ \Delta^{(\ell)}_N $ is
important to know the reducibility of $V^{\delta,p}_{\ell}$.
\begin{proposition}\label{prop1}
The Shapovalov determinants $ \Delta^{(\ell)}_N$ at level~$N$ of ${\mathfrak g}_{\ell}$ are given as follows $($up to
overall sign$)$:
\begin{alignat}{3}
   & i)
\quad &&
\Delta_N^{(1)} = \left(\prod\limits_{m=0}^N m! \right)^2 (2p)^{N(N+1)},&
\label{KDl1}
\\
  &   ii)
\quad &&
\Delta^{(\ell)}_N = (\ell+1)^2   p^2   \delta_{N1},
\qquad
\ell \geq 2.&
%\label{KDell}
\end{alignat}
\end{proposition}

\begin{proof}
The case with $ \ell = 1 $ shows a~deviation from other values of $ \ell$.
We treat the case with $ \ell = 1 $ separately.
The basis of $ V^{\delta,p}_1 $ is specif\/ied by two nonnegative integers:
\begin{gather*}
\ket{k,m} = H^k P_0^m \ket{\delta,p}.
\end{gather*}
The next relation obtained from~\eqref{PnHkcomm} is useful in the following computation:
\begin{gather}
P_2 \ket{k,m} = 2kp \ket{k-1,m} + k(k-1) \ket{k-2,m+1}.
\label{ActionP2}
\end{gather}
The level is $ N = k +m $ so that the basis of $ \big(V^{\delta,p}_1\big)_N $ is given~by
\begin{gather*}
\ket{m,N-m},
\qquad
0 \leq m \leq N.
\end{gather*}
The product of two vectors in $ \big(V^{\delta,p}_1\big)_N $ is
\begin{gather}
\bracket{N-k,k}{m,N-m} = \bra{N-k,0} P_0^{N-m} P_2^k \ket{m,0}.
\label{IPl=1}
\end{gather}

It follows from~\eqref{ActionP2} that $P_2^k \ket{m,0}$ is a~linear combination of $ \ket{m-k-j,j} $ with $0 \leq j
\leq k$.
Therefore, $ P_2^k \ket{m,0} = 0$ if $k > m$.
Thus we have proved the following lemma:

\begin{lemma}
\label{lemma1}
$ \bracket{N-k,k}{m,N-m} = 0 $ if $ k > m$.
\end{lemma}

By def\/inition $ \Delta^{(1)}_N $ is given by (up to sign)
\begin{gather*}
\Delta^{(1)}_N = \left|
\begin{matrix}
 \bracket{N,0}{0,N} & \bracket{N,0}{1,N-1} & \dots & \bracket{N,0}{N,0}
\\
\bracket{N-1,0}{0,N} & \bracket{N-1,0}{1,N-1} & \dots & \bracket{N-1,0}{N,0}
\\
\vdots & \vdots & \ddots & \vdots
\\
\bracket{0,N}{0,N} & \bracket{0,N}{1,N-1} & \dots & \bracket{0N}{N,0}
\end{matrix}
\right|.
\end{gather*}
By Lemma~\ref{lemma1}, $ \Delta^{(1)}_N $ is the determinant of upper triangular matrix.
Thus
\begin{gather*}
\Delta_N^{(1)} = \prod\limits_{m=0}^N \bracket{N-m,m}{m,N-m}.
\end{gather*}
Each factor is calculated by~\eqref{IPl=1} as follows
\begin{gather*}
\bracket{N-m,m}{m,N-m} = \bra{N-m,m} P_0^{N-m} P_2^m \ket{N-m,m} = (N-m)!   m!   (2p)^N,
\end{gather*}
where we used the relation obtained from~\eqref{ActionP2}: $ P_2^m \ket{m,0} = (2p)^m m ! \ket{0,0}$.
This completes the proof of~\eqref{KDl1}.

Now let us turn to the case with $\ell \geq 2$.
In this case $ \Delta^{(\ell)}_N \neq 0 $ only if $N = 1$.
As we shall see, this fact stems from that at least two rows of $ \Delta^{(\ell)}_N $ are proportional if $ N > 1$.
The basis of $(V^{\delta,p}_{\ell})_N $ is given by~\eqref{anyl-basis} with the constraint $  N = k+
\sum\limits_{i=1}^{\ell} i   m_i$.
First we prove the following lemma:
\begin{lemma}
\label{lemma2}
In the subspace $ (V^{\delta,p}_{\ell})_N $ we have the relation
\begin{gather*}
\bracket{0,\underaccent{\tilde}{n}}{k,\underaccent{\tilde}{m}} = f^{(\ell)}_{N,\underaccent{\tilde}{n}}(p) \delta_{kN},
%\label{IPlge2}
\end{gather*}
where $ f^{(\ell)}_{N,\underaccent{\tilde}{n}}(p) $ is a~function of~$p$ determined by $ \ell$, $N $ and $
\ket{0,\underaccent{\tilde}{n}}$.
\end{lemma}
\begin{proof}
By def\/inition
\begin{gather}
    \bracket{0,\underaccent{\tilde}{n}}{k,\underaccent{\tilde}{m}} = \bra{\delta,p} P_{2\ell}^{n_{\ell}}
P_{2\ell-1}^{n_{\ell-1}} \cdots P_{\ell+1}^{n_1} H^k P_{\ell-1}^{m_1} \cdots P_0^{m_{\ell}} \ket{\delta,p},
\label{IPlge2-2}
\\
    N = \sum\limits_{i=1}^{\ell} i n_i = k + \sum\limits_{i=1}^{\ell} i  m_i.
\nonumber
\end{gather}
Because of~\eqref{PnHkcomm} the right hand side of~\eqref{IPlge2-2} will be a~linear combination of the terms
\begin{gather*}
\bra{\delta,p} H^{k'} P_{\ell-1}^{m'_1} \cdots P_0^{m'_{\ell}} \ket{\delta,p},
\end{gather*}
with the condition $ m'_k \geq m_k $ for all~$k$.
However, such terms give nonvanishing contributions only if
\begin{gather*}
k' = m'_1 = m'_2 = \dots = m'_{\ell} = 0.
\end{gather*}
This implies $ m_k = 0 $ for all $k$, i.e., $ k = N$.
\end{proof}

It follows from Lemma~\ref{lemma2} that two rows in $ \Delta^{(\ell)}_N $ labelled by $ \bra{0,\underaccent{\tilde}{n}}
$ and $ \bra{0,\underaccent{\tilde}{n'}} $ are proportional.
If $ N \geq 2 $ then there exists at least one such pair of rows, but no such pair for $N = 1$.
Thus $ \Delta^{(\ell)}_N = 0 $ if $ N \geq 2$.
On the other hand the $ N = 1 $ subspace is two-dimensional with the basis $ \ket{1, \underaccent{\tilde}{0}}, \ket{0,
\underaccent{\tilde}{\epsilon}_1} $ where $\underaccent{\tilde}{0}$ denotes the zero vector in ${\mathbb R}^{\ell}$.
Thus $ \Delta^{(\ell)}_1 $ is computed as follows
\begin{gather*}
\Delta^{(\ell)}_1 = \left|
\begin{matrix} \bracket{1, \underaccent{\tilde}{0}}{1, \underaccent{\tilde}{0}} & \bracket{1, \underaccent{\tilde}{0}}{0,
\underaccent{\tilde}{\epsilon}_1}
\vspace{1mm}\\  \bracket{0, \underaccent{\tilde}{\epsilon}_1}{1, \underaccent{\tilde}{0}} & \bracket{0,
\underaccent{\tilde}{\epsilon}_1}{0, \underaccent{\tilde}{\epsilon}_1}
\end{matrix}
\right| = \left|
\begin{matrix} 2 \delta & (\ell+1)p
\vspace{1mm}\\  (\ell+1)p & 0
\end{matrix}
\right| = - (\ell+1)^2   p^2   .
\end{gather*}
This completes the proof of Proposition~\ref{prop1}.
\end{proof}



\begin{proposition}\qquad
\label{prop2}
\begin{enumerate}\itemsep=0pt
\item[$i)$] $ V^{\delta,p}_{\ell} $ is irreducible if $ \ell = 1 $ and $ p \neq 0$.
\item[$ii)$] $ V^{\delta,p}_{\ell} $ is reducible if $ \ell \geq 2 $ and $ p \neq 0$.
\item[$iii)$] $ V^{\delta,0}_{\ell} $ is reducible for all values of $ \ell$.
\end{enumerate}
\end{proposition}



\begin{proof}
(i) is the corollary of Proposition~\ref{prop1}(i).
Proposition~\ref{prop1} also suggests the existence of singular vectors in $V^{\delta,0}_{\ell} $ for any $\ell$ and
in $V^{\delta,p}_{\ell}$ with $\ell \geq 2$ and $p \neq 0$.
In the next section we shall show that this is indeed the case.
Thus we establish~(ii) and~(iii).
\end{proof}
